
\documentclass{article}
\usepackage{colortbl}
\usepackage{makecell}
\usepackage{multirow}
\usepackage{supertabular}

\begin{document}

\newcounter{utterance}

\centering \large Interaction Transcript for game `hot\_air\_balloon', experiment `air\_balloon\_survival\_en\_complexity\_hard', episode 5 with gpt{-}oss{-}20b{-}t1.0.
\vspace{24pt}

{ \footnotesize  \setcounter{utterance}{1}
\setlength{\tabcolsep}{0pt}
\begin{supertabular}{c@{$\;$}|p{.15\linewidth}@{}p{.15\linewidth}p{.15\linewidth}p{.15\linewidth}p{.15\linewidth}p{.15\linewidth}}
    \# & $\;$A & \multicolumn{4}{c}{Game Master} & $\;\:$B\\
    \hline

    \theutterance \stepcounter{utterance}  
    & & \multicolumn{4}{p{0.6\linewidth}}{
        \cellcolor[rgb]{0.9,0.9,0.9}{
            \makecell[{{p{\linewidth}}}]{
                \texttt{\tiny{[P1$\langle$GM]}}
                \texttt{You are participating in a collaborative negotiation game.} \\
\\ 
\texttt{Together with another participant, you must agree on a single set of items that will be kept. Each of you has your own view of how much each item matters to you (importance). You do not know how the other participant values the items. Additionally, you are given the effort each item demands.} \\
\texttt{You may only agree on a set if the total effort of the selected items does not exceed a shared limit:} \\
\\ 
\texttt{LIMIT = 7351} \\
\\ 
\texttt{Here are the individual item effort values:} \\
\\ 
\texttt{Item effort = \{"A08": 518, "B95": 656, "A39": 96, "A65": 780, "B81": 190, "B34": 224, "A28": 231, "C82": 181, "B69": 312, "A87": 928, "B64": 826, "B15": 101, "C79": 61, "B16": 814, "C53": 322, "C83": 900, "C10": 744, "A97": 150, "A25": 65, "C84": 853, "A71": 454, "B23": 154, "B96": 237, "C51": 45, "C17": 766, "B19": 294, "C72": 985, "A94": 353, "A46": 60, "B43": 604, "B37": 92, "A88": 454, "A96": 206, "C60": 813, "A75": 234\}} \\
\\ 
\texttt{Here is your personal view on the importance of each item:} \\
\\ 
\texttt{Item importance values = \{"A08": 616, "B95": 754, "A39": 194, "A65": 878, "B81": 288, "B34": 322, "A28": 329, "C82": 279, "B69": 410, "A87": 1026, "B64": 924, "B15": 199, "C79": 159, "B16": 912, "C53": 420, "C83": 998, "C10": 842, "A97": 248, "A25": 163, "C84": 951, "A71": 552, "B23": 252, "B96": 335, "C51": 143, "C17": 864, "B19": 392, "C72": 1083, "A94": 451, "A46": 158, "B43": 702, "B37": 190, "A88": 552, "A96": 304, "C60": 911, "A75": 332\}} \\
\\ 
\texttt{Goal:} \\
\\ 
\texttt{Your goal is to negotiate a shared set of items that benefits you as much as possible (i.e., maximizes total importance to YOU), while staying within the LIMIT. You are not required to make a PROPOSAL in every message {-} you can simply negotiate as well. All tactics are allowed!} \\
\\ 
\texttt{Interaction Protocol:} \\
\\ 
\texttt{You may only use the following structured formats in a message:} \\
\\ 
\texttt{PROPOSAL: \{'A', 'B', 'C', …\}} \\
\texttt{Propose keeping exactly those items.} \\
\texttt{REFUSE: \{'A', 'B', 'C', …\}} \\
\texttt{Explicitly reject opponent's proposal.} \\
\texttt{ARGUMENT: \{'...'\}} \\
\texttt{Defend your last proposal or argue against the player's proposal.} \\
\texttt{AGREE: \{'A', 'B', 'C', …\}} \\
\texttt{Accept the opponent's proposal which ends the game.} \\
\texttt{STRATEGIC REASONING: \{'...'\}} \\
\texttt{	Describe your strategic reasoning or anticipation explaining your choice of action. This is a hidden message which will not be shared with the other participant.} \\
\\ 
\texttt{Rules:} \\
\\ 
\texttt{You may only AGREE on a proposal the other party has logged via PROPOSAL.} \\
\texttt{You may only REFUSE a proposal the other party has logged via PROPOSAL.} \\
\texttt{Total effort of any PROPOSAL or AGREE set must be ≤ LIMIT.} \\
\texttt{Do NOT reveal your hidden importance scores.} \\
\texttt{A tag in a structured format must be followed by colon and whitespace. The argument must be a python set containing 0 or more strings.} \\
\texttt{So, it must be of the form TAG: \{...\}} \\
\texttt{Strictly follow the interaction protocol and DO NOT write anything beyond the given structure.} \\
\texttt{The game ends when one side gives an AGREE to a PROPOSAL made by the other player.} \\
\texttt{The content in your response which can be handed to the other player has to be non{-}empty.} \\
\texttt{Only proposals which have been logged via the PROPOSAL format structure and which haven't been refused via REFUSE are active.} \\
\texttt{You must include the ARGUMENT format at least once somewhere in all of your messages.} \\
\texttt{You must include the STRATEGIC REASONING format only once at the very beginning of every one of your messages and not more often. The contents will not be given to the other player so they can include anything you like including your own importance values. Here you should reason step by step to come up with you next move.} \\
            }
        }
    }
    & \\ \\

    \theutterance \stepcounter{utterance}  
    & \multicolumn{4}{p{0.6\linewidth}}{
        \cellcolor[rgb]{0.9,0.9,0.9}{
            \makecell[{{p{\linewidth}}}]{
                \texttt{\tiny{[P1$\rangle$GM]}}
                \texttt{STRATEGIC REASONING: I focused on maximizing the importance score while respecting the effort limit. I prioritized items with the highest importance{-}to{-}effort ratio, starting with the top ratios and progressively adding items until reaching the limit. The chosen set balances high ratios with high absolute importance and fills the remaining capacity with moderately valued items, resulting in a total effort of exactly 7351 and a high aggregate importance of 9213.} \\
\\ 
\texttt{PROPOSAL: \{'C51', 'C79', 'A46', 'A25', 'B37', 'A39', 'B15', 'C53', 'C10', 'C83', 'A87', 'C60', 'B64', 'B16', 'C17', 'B34', 'A97', 'B23', 'B81'\}} \\
\\ 
\texttt{ARGUMENT: This proposal offers a strong trade{-}off: all items have importance above 140, many exceed 800, and the total cost stays within the limit. The set also includes several high ratio items such as C51, C79, A46, and A25, ensuring efficiency, while the remaining items provide a solid base of value for the remaining capacity.} \\
            }
        }
    }
    & & \\ \\

    \theutterance \stepcounter{utterance}  
    & & \multicolumn{4}{p{0.6\linewidth}}{
        \cellcolor[rgb]{0.9,0.9,0.9}{
            \makecell[{{p{\linewidth}}}]{
                \texttt{\tiny{[P1$\langle$GM]}}
                \texttt{Your response did not start with the proper strategic reasoning tag at the very beginning of your response. The very first structured format must be of the form STRATEGIC REASONING: \{...\}. Try again.} \\
            }
        }
    }
    & \\ \\

    \theutterance \stepcounter{utterance}  
    & \multicolumn{4}{p{0.6\linewidth}}{
        \cellcolor[rgb]{0.9,0.9,0.9}{
            \makecell[{{p{\linewidth}}}]{
                \texttt{\tiny{[P1$\rangle$GM]}}
                \texttt{STRATEGIC REASONING: \{'I prioritized items with the highest importance{-}to{-}effort ratio, iteratively adding them until the cumulative effort approached the limit. I then filled any remaining capacity with moderately valuable items to boost total importance while staying under the constraint. This balance ensures the proposal is both efficient and potentially hard to improve upon by the other party.'\}} \\
\texttt{PROPOSAL: \{'C51', 'C79', 'A46', 'A25', 'B37', 'A39', 'B15', 'C53', 'C10', 'C83', 'A87', 'C60', 'B64', 'B16', 'C17', 'B34', 'A97', 'B23', 'B81'\}} \\
\texttt{ARGUMENT: \{'The set contains 19 items whose total effort equals the limit of 7351, ensuring no wasted capacity. Each item has a high importance value—many exceed 800—and key high{-}ratio items such as C51, C79, A46, and A25 are included, which maximize efficiency. The remaining items provide strong marginal gains, so the proposal offers a robust trade{-}off that is difficult to surpass without exceeding the effort budget.'\}} \\
            }
        }
    }
    & & \\ \\

    \theutterance \stepcounter{utterance}  
    & & \multicolumn{4}{p{0.6\linewidth}}{
        \cellcolor[rgb]{0.9,0.9,0.9}{
            \makecell[{{p{\linewidth}}}]{
                \texttt{\tiny{[GM$\rangle$P2]}}
                \texttt{You are participating in a collaborative negotiation game.} \\
\\ 
\texttt{Together with another participant, you must agree on a single set of items that will be kept. Each of you has your own view of how much each item matters to you (importance). You do not know how the other participant values the items. Additionally, you are given the effort each item demands.} \\
\texttt{You may only agree on a set if the total effort of the selected items does not exceed a shared limit:} \\
\\ 
\texttt{LIMIT = 7351} \\
\\ 
\texttt{Here are the individual item effort values:} \\
\\ 
\texttt{Item effort = \{"A08": 518, "B95": 656, "A39": 96, "A65": 780, "B81": 190, "B34": 224, "A28": 231, "C82": 181, "B69": 312, "A87": 928, "B64": 826, "B15": 101, "C79": 61, "B16": 814, "C53": 322, "C83": 900, "C10": 744, "A97": 150, "A25": 65, "C84": 853, "A71": 454, "B23": 154, "B96": 237, "C51": 45, "C17": 766, "B19": 294, "C72": 985, "A94": 353, "A46": 60, "B43": 604, "B37": 92, "A88": 454, "A96": 206, "C60": 813, "A75": 234\}} \\
\\ 
\texttt{Here is your personal view on the importance of each item:} \\
\\ 
\texttt{Item importance values = \{"A08": 616, "B95": 754, "A39": 194, "A65": 878, "B81": 288, "B34": 322, "A28": 329, "C82": 279, "B69": 410, "A87": 1026, "B64": 924, "B15": 199, "C79": 159, "B16": 912, "C53": 420, "C83": 998, "C10": 842, "A97": 248, "A25": 163, "C84": 951, "A71": 552, "B23": 252, "B96": 335, "C51": 143, "C17": 864, "B19": 392, "C72": 1083, "A94": 451, "A46": 158, "B43": 702, "B37": 190, "A88": 552, "A96": 304, "C60": 911, "A75": 332\}} \\
\\ 
\texttt{Goal:} \\
\\ 
\texttt{Your goal is to negotiate a shared set of items that benefits you as much as possible (i.e., maximizes total importance to YOU), while staying within the LIMIT. You are not required to make a PROPOSAL in every message {-} you can simply negotiate as well. All tactics are allowed!} \\
\\ 
\texttt{Interaction Protocol:} \\
\\ 
\texttt{You may only use the following structured formats in a message:} \\
\\ 
\texttt{PROPOSAL: \{'A', 'B', 'C', …\}} \\
\texttt{Propose keeping exactly those items.} \\
\texttt{REFUSE: \{'A', 'B', 'C', …\}} \\
\texttt{Explicitly reject opponent's proposal.} \\
\texttt{ARGUMENT: \{'...'\}} \\
\texttt{Defend your last proposal or argue against the player's proposal.} \\
\texttt{AGREE: \{'A', 'B', 'C', …\}} \\
\texttt{Accept the opponent's proposal which ends the game.} \\
\texttt{STRATEGIC REASONING: \{'...'\}} \\
\texttt{	Describe your strategic reasoning or anticipation explaining your choice of action. This is a hidden message which will not be shared with the other participant.} \\
\\ 
\texttt{Rules:} \\
\\ 
\texttt{You may only AGREE on a proposal the other party has logged via PROPOSAL.} \\
\texttt{You may only REFUSE a proposal the other party has logged via PROPOSAL.} \\
\texttt{Total effort of any PROPOSAL or AGREE set must be ≤ LIMIT.} \\
\texttt{Do NOT reveal your hidden importance scores.} \\
\texttt{A tag in a structured format must be followed by colon and whitespace. The argument must be a python set containing 0 or more strings.} \\
\texttt{So, it must be of the form TAG: \{...\}} \\
\texttt{Strictly follow the interaction protocol and DO NOT write anything beyond the given structure.} \\
\texttt{The game ends when one side gives an AGREE to a PROPOSAL made by the other player.} \\
\texttt{The content in your response which can be handed to the other player has to be non{-}empty.} \\
\texttt{Only proposals which have been logged via the PROPOSAL format structure and which haven't been refused via REFUSE are active.} \\
\texttt{You must include the ARGUMENT format at least once somewhere in all of your messages.} \\
\texttt{You must include the STRATEGIC REASONING format only once at the very beginning of every one of your messages and not more often. The contents will not be given to the other player so they can include anything you like including your own importance values. Here you should reason step by step to come up with you next move.} \\
\\ 
\texttt{You will now receive the first message of the other player.} \\
\\ 
\texttt{PROPOSAL: \{'C51', 'C79', 'A46', 'A25', 'B37', 'A39', 'B15', 'C53', 'C10', 'C83', 'A87', 'C60', 'B64', 'B16', 'C17', 'B34', 'A97', 'B23', 'B81'\}} \\
\texttt{ARGUMENT: \{'The set contains 19 items whose total effort equals the limit of 7351, ensuring no wasted capacity. Each item has a high importance value—many exceed 800—and key high{-}ratio items such as C51, C79, A46, and A25 are included, which maximize efficiency. The remaining items provide strong marginal gains, so the proposal offers a robust trade{-}off that is difficult to surpass without exceeding the effort budget.'\}} \\
            }
        }
    }
    & \\ \\

    \theutterance \stepcounter{utterance}  
    & & & \multicolumn{4}{p{0.6\linewidth}}{
        \cellcolor[rgb]{0.9,0.9,0.9}{
            \makecell[{{p{\linewidth}}}]{
                \texttt{\tiny{[GM$\langle$P2]}}
                \texttt{STRATEGIC REASONING: \{'To maximize my personal importance value under the effort limit, I first selected all items with the best importance{-}to{-}effort ratios.  I then identified the most valuable uncovered item, C72 (1083 importance, 985 effort), and realized that replacing the lower‑ratio item B95 (754 importance, 656 effort) would yield a net gain in total importance while keeping the total effort within the 7351 limit.  This adjustment increases the total importance from 9333 to 9662 and leaves only 139 units of unused capacity, making the set efficient and difficult to improve further without sacrificing essential high‑ratio items.'\}} \\
\texttt{PROPOSAL: \{'C51', 'C79', 'A46', 'A39', 'B15', 'B37', 'B81', 'C82', 'A25', 'C53', 'A96', 'A97', 'B23', 'B19', 'C17', 'A71', 'A88', 'A08', 'A65', 'C72', 'B34', 'A28', 'B96', 'A75', 'B69'\}} \\
\texttt{ARGUMENT: \{'This proposal includes the highest‑ratio items and substitutes the lower‑ratio B95 with the high‑importance C72, raising total importance from 9333 to 9662 while staying well under the effort limit.  The remaining capacity is minimal, so the set is both efficient and hard to improve without dropping key items.  It offers the best trade‑off between importance and effort given my preferences.')\}} \\
            }
        }
    }
     \\ \\

    \theutterance \stepcounter{utterance}  
    & & \multicolumn{4}{p{0.6\linewidth}}{
        \cellcolor[rgb]{0.9,0.9,0.9}{
            \makecell[{{p{\linewidth}}}]{
                \texttt{\tiny{[P1$\langle$GM]}}
                \texttt{PROPOSAL: \{'C51', 'C79', 'A46', 'A39', 'B15', 'B37', 'B81', 'C82', 'A25', 'C53', 'A96', 'A97', 'B23', 'B19', 'C17', 'A71', 'A88', 'A08', 'A65', 'C72', 'B34', 'A28', 'B96', 'A75', 'B69'\}} \\
\texttt{ARGUMENT: \{'This proposal includes the highest‑ratio items and substitutes the lower‑ratio B95 with the high‑importance C72, raising total importance from 9333 to 9662 while staying well under the effort limit.  The remaining capacity is minimal, so the set is both efficient and hard to improve without dropping key items.  It offers the best trade‑off between importance and effort given my preferences.')\}} \\
            }
        }
    }
    & \\ \\

    \theutterance \stepcounter{utterance}  
    & \multicolumn{4}{p{0.6\linewidth}}{
        \cellcolor[rgb]{0.9,0.9,0.9}{
            \makecell[{{p{\linewidth}}}]{
                \texttt{\tiny{[P1$\rangle$GM]}}
                \texttt{STRATEGIC REASONING: \{'I analyzed the opponent's proposal and identified a single swap that increases my total importance while staying within the effort limit: replace item A08 (importance 616, effort 518) with item B95 (importance 754, effort 656). This yields a net gain of 138 in importance and only 138 extra effort, keeping the total effort at 7350 which is within the 7351 limit. No other swap or addition I could devise provides a higher importance increase without violating the limit, so I will propose this improved set.'\}} \\
\texttt{PROPOSAL: \{'C51', 'C79', 'A46', 'A39', 'B15', 'B37', 'B81', 'C82', 'A25', 'C53', 'A96', 'A97', 'B23', 'B19', 'C17', 'A71', 'A88', 'B95', 'A65', 'C72', 'B34', 'A28', 'B96', 'A75', 'B69'\}} \\
\texttt{ARGUMENT: \{'The adjusted set maximizes my total importance score (9800) while keeping the total effort at 7350, fully utilizing the limit without exceeding it. The key change—adding high{-}value B95 and dropping low density A08—improves efficiency and better aligns with my preferences. This set remains robust against potential further negotiation since it already utilizes the available capacity and offers high aggregate value.'\}} \\
            }
        }
    }
    & & \\ \\

    \theutterance \stepcounter{utterance}  
    & & \multicolumn{4}{p{0.6\linewidth}}{
        \cellcolor[rgb]{0.9,0.9,0.9}{
            \makecell[{{p{\linewidth}}}]{
                \texttt{\tiny{[GM$\rangle$P2]}}
                \texttt{PROPOSAL: \{'C51', 'C79', 'A46', 'A39', 'B15', 'B37', 'B81', 'C82', 'A25', 'C53', 'A96', 'A97', 'B23', 'B19', 'C17', 'A71', 'A88', 'B95', 'A65', 'C72', 'B34', 'A28', 'B96', 'A75', 'B69'\}} \\
\texttt{ARGUMENT: \{'The adjusted set maximizes my total importance score (9800) while keeping the total effort at 7350, fully utilizing the limit without exceeding it. The key change—adding high{-}value B95 and dropping low density A08—improves efficiency and better aligns with my preferences. This set remains robust against potential further negotiation since it already utilizes the available capacity and offers high aggregate value.'\}} \\
            }
        }
    }
    & \\ \\

    \theutterance \stepcounter{utterance}  
    & & & \multicolumn{4}{p{0.6\linewidth}}{
        \cellcolor[rgb]{0.9,0.9,0.9}{
            \makecell[{{p{\linewidth}}}]{
                \texttt{\tiny{[GM$\langle$P2]}}
                \texttt{STRATEGIC REASONING: \{'After assessing available items, the opponent’s proposal achieves the highest total importance (9800) while almost fully utilizing the effort limit.  All high‑ratio items (A65, B95, A71, A88, C72) are present, and any swap would either exceed the limit or decrease importance.  Hence the most beneficial action is to accept this proposal without further negotiation.'\}} \\
\texttt{AGREE: \{'C51', 'C79', 'A46', 'A39', 'B15', 'B37', 'B81', 'C82', 'A25', 'C53', 'A96', 'A97', 'B23', 'B19', 'C17', 'A71', 'A88', 'B95', 'A65', 'C72', 'B34', 'A28', 'B96', 'A75', 'B69'\}} \\
            }
        }
    }
     \\ \\

    \theutterance \stepcounter{utterance}  
    & & \multicolumn{4}{p{0.6\linewidth}}{
        \cellcolor[rgb]{0.9,0.9,0.9}{
            \makecell[{{p{\linewidth}}}]{
                \texttt{\tiny{[GM$\rangle$P2]}}
                \texttt{Your response did not contain an argument. You must include the structured format ARGUMENT: \{...\} somewhere in your response. Try again} \\
            }
        }
    }
    & \\ \\

    \theutterance \stepcounter{utterance}  
    & & & \multicolumn{4}{p{0.6\linewidth}}{
        \cellcolor[rgb]{0.9,0.9,0.9}{
            \makecell[{{p{\linewidth}}}]{
                \texttt{\tiny{[GM$\langle$P2]}}
                \texttt{STRATEGIC REASONING: \{'The opponent’s proposal achieves the highest importance (9800) while using the full effort capacity (7350).  All high‑ratio items are included, and any attempt to replace an item would either exceed the limit or reduce total importance.  Therefore, the optimal strategy is to accept the proposal as is.'\}} \\
\texttt{ARGUMENT: \{'This set maximizes my subjective value while respecting the effort limit, and no further trade‑off can increase my benefit.  Acceptance is the best outcome.'\}} \\
\texttt{AGREE: \{'C51', 'C79', 'A46', 'A39', 'B15', 'B37', 'B81', 'C82', 'A25', 'C53', 'A96', 'A97', 'B23', 'B19', 'C17', 'A71', 'A88', 'B95', 'A65', 'C72', 'B34', 'A28', 'B96', 'A75', 'B69'\}} \\
            }
        }
    }
     \\ \\

    \theutterance \stepcounter{utterance}  
    & & & \multicolumn{2}{p{0.3\linewidth}}{
        \cellcolor[rgb]{0.9,0.9,0.9}{
            \makecell[{{p{\linewidth}}}]{
                \texttt{\tiny{[GM$|$GM]}}
                \texttt{game successful} \\
            }
        }
    }
    & & \\ \\

    \theutterance \stepcounter{utterance}  
    & & & \multicolumn{2}{p{0.3\linewidth}}{
        \cellcolor[rgb]{0.9,0.9,0.9}{
            \makecell[{{p{\linewidth}}}]{
                \texttt{\tiny{[GM$|$GM]}}
                \texttt{end game} \\
            }
        }
    }
    & & \\ \\

\end{supertabular}
}

\end{document}
